\section{The actual tangent germ residue field}

Write $B=\mathcal O_{\mathbb C^d,0}$, $\mathfrak m_B$ for its genuine
maximal ideal, and $\kappa(B)=B/\mathfrak m_B$. Complex coefficients carry
the actual $B$-algebra structure given by evaluation at the origin.

\begin{theorem}[Canonical residue field and literal origin values]
\label{hir:actual-base-residue-field}
\lean{Hironaka.HironakaActualBaseResidueField.baseResidueFieldEquiv,Hironaka.HironakaActualBaseResidueField.baseResidueFieldEquiv_residue,Hironaka.HironakaActualBaseResidueField.baseResidueFieldEquiv_symm}
\leanok
Evaluation induces a $B$-algebra isomorphism
$e:\kappa(B)\simeq_B\mathbb C$ satisfying $e([g])=g(0)$ for every
original holomorphic germ $g$. Its inverse sends a complex number to the
class of the corresponding actual constant germ.
\end{theorem}
\begin{proof}
The full maximal ideal consists exactly of the germs whose value at the
origin is zero. Evaluation therefore descends to the quotient. Its map
is injective because its domain is a field and the map preserves one.
Every complex number is the value of its constant germ, proving
surjectivity and the stated inverse formula. The formula on a residue
class proves compatibility with the original $B$-algebra structures.
\end{proof}

\begin{theorem}[Natural tensor comparison and detection of injectivity]
\label{hir:actual-base-residue-tensor}
\lean{Hironaka.HironakaActualBaseResidueField.residueTensorEquiv,Hironaka.HironakaActualBaseResidueField.residueTensorEquiv_tmul,Hironaka.HironakaActualBaseResidueField.residueTensorEquiv_natural,Hironaka.HironakaActualBaseResidueField.residue_lTensor_injective_iff_complex}
\leanok
\uses{hir:actual-base-residue-field}
For every $B$-module $P$, the formula $a\otimes p\mapsto e(a)\otimes p$
gives a $B$-linear equivalence $E_P:\kappa(B)\otimes_B P\simeq_B
\mathbb C\otimes_B P$. For every $B$-linear map $f:P\to Q$ it satisfies
$E_Q(1\otimes f)=(1\otimes f)E_P$. Consequently tensoring $f$ with
$\kappa(B)$ is injective exactly when tensoring it with the actual
complex residue algebra is injective.
\end{theorem}
\begin{proof}
Tensor the underlying $B$-linear equivalence of $e$ with the identity of
$P$. Its inverse tensors $e^{-1}$ with that identity. Both sides of the
naturality equality send $a\otimes p$ to $e(a)\otimes f(p)$; additivity
extends the equality to every tensor. The two vertical maps of this
commuting square are bijective, so injectivity of either horizontal
map is equivalent to that of the other.
\end{proof}

\begin{theorem}[The complex residue criterion for actual filtered cuts]
\label{hir:actual-base-complex-symbol-criterion}
\lean{Hironaka.HironakaActualBaseResidueField.adicLayer_after_parameter_finite_free_of_complex}
\leanok
\uses{hir:actual-base-residue-tensor,hir:local-actual-symbol-quotient-layers-free}
Let $R$ be a commutative $B$-algebra, $I,N\subset R$ ideals, and
$v\in N$. Suppose every actual adic layer
$G_q=\operatorname{im}(N^q\to R/(I+N^{q+1}))$ is finite free over $B$,
and every literal multiplication map $G_q\to G_{q+1}$ becomes injective
under the actual complex residue tensor. Then all actual layers for
$I+(v)$ and the same $N$ are finite free over $B$.
\end{theorem}
\begin{proof}
Apply the tensor comparison to every actual multiplication map,
restricted to $B$. It supplies exactly the canonical residue-field
injectivity hypothesis of the existing filtered-cut theorem. That
theorem splits the multiplication map, proves the whole filtered
intersection equality, and identifies the next actual layer with its
free quotient. Its degree-zero argument uses the genuine zero-degree
layer isomorphism; no positive-degree assumption is imposed here.
\end{proof}
